% Source: https://artofproblemsolving.com/community/c3h3248681_202324_abmc_online_december_contest__actonboxborough_math_competition
% Stable identity: aops-topic-3248681-q10
Place $A=(0,0),B=(51,0)$. The distance equations give $C=(68,68\sqrt3)$, since $68^2+(68\sqrt3)^2=136^2$ and $17^2+(68\sqrt3)^2=119^2$. Thus $C^{\prime}=(68,-68\sqrt3)$ and line $AC^{\prime}$ makes angle $-60^\circ$ with the positive horizontal axis. Reflecting $B$ across it yields $B^{\prime}=51(\cos(-120^\circ),\sin(-120^\circ))=(-\dfrac{51}{2},-\dfrac{51\sqrt3}{2})$. The vertical base $CC^{\prime}$ has length $136\sqrt3$ and the perpendicular height from $B^{\prime}$ is $68+\dfrac{51}{2}=\dfrac{187}{2}$. The area is $\dfrac{1}{2}(136\sqrt3)\dfrac{187}{2}=6358\sqrt3$, giving $\boxed{a+b=6361}$.
\begin{asy}
usepackage("amsmath");
size(7cm);
pair A=(0,0),B=(51,0),C=(68,68sqrt(3)),Cp=(68,-68sqrt(3)),Bp=(-25.5,-25.5sqrt(3));draw(A--B--C--cycle);draw(A--Cp,dashed);draw(Bp--Cp--C--cycle,blue);draw(B--Bp,dashed);label("$A$",A,W);label("$B$",B,N);label("$C$",C,N);label("$C^{\prime}$",Cp,S);label("$B^{\prime}$",Bp,W);
\end{asy}
